\documentclass[12pt, a4paper]{article}

% ── Paquetes ─────────────────────────────────────────────────────────
\usepackage[utf8]{inputenc}
\usepackage[T1]{fontenc}
% babel-spanish no disponible en esta instalación
\usepackage{amsmath, amssymb, amsthm}
\usepackage{geometry}
\usepackage{graphicx}
\usepackage{xcolor}
\usepackage{mdframed}
\usepackage{enumitem}
\usepackage{booktabs}
\usepackage{hyperref}

% ── Geometría ────────────────────────────────────────────────────────
\geometry{
    top    = 2.5cm,
    bottom = 2.5cm,
    left   = 3.0cm,
    right  = 3.0cm
}

% ── Colores ──────────────────────────────────────────────────────────
\definecolor{azulviz}{RGB}{37, 99, 235}
\definecolor{naranjviz}{RGB}{251, 146, 60}
\definecolor{verdviz}{RGB}{52, 211, 153}
\definecolor{amarillviz}{RGB}{251, 191, 36}
\definecolor{fondocaja}{RGB}{240, 245, 255}
\definecolor{bordecruz}{RGB}{37, 99, 235}
\definecolor{fondoadv}{RGB}{255, 248, 230}
\definecolor{bordeadv}{RGB}{251, 146, 60}
\definecolor{fondoteo}{RGB}{235, 250, 242}
\definecolor{bordeteo}{RGB}{52, 153, 110}

% ── Entornos ─────────────────────────────────────────────────────────
\mdfdefinestyle{vizverde}{
    backgroundcolor = fondoteo,
    linecolor       = bordeteo,
    linewidth       = 1.5pt,
    leftmargin      = 0pt,
    rightmargin     = 0pt,
    innertopmargin  = 8pt,
    innerbottommargin = 8pt,
    innerleftmargin  = 12pt,
    innerrightmargin = 12pt,
    roundcorner     = 4pt
}

\mdfdefinestyle{viznaranja}{
    backgroundcolor = fondoadv,
    linecolor       = bordeadv,
    linewidth       = 1.5pt,
    leftmargin      = 0pt,
    rightmargin     = 0pt,
    innertopmargin  = 8pt,
    innerbottommargin = 8pt,
    innerleftmargin  = 12pt,
    innerrightmargin = 12pt,
    roundcorner     = 4pt
}

\mdfdefinestyle{vizazul}{
    backgroundcolor = fondocaja,
    linecolor       = bordecruz,
    linewidth       = 1.5pt,
    leftmargin      = 0pt,
    rightmargin     = 0pt,
    innertopmargin  = 8pt,
    innerbottommargin = 8pt,
    innerleftmargin  = 12pt,
    innerrightmargin = 12pt,
    roundcorner     = 4pt
}

\newenvironment{viznota}[1]{%
    \begin{mdframed}[style=vizverde]
    \textbf{\color{bordeteo}En GradienViz \textemdash\ #1.}\quad
}{%
    \end{mdframed}
}

\newenvironment{vizadv}{%
    \begin{mdframed}[style=viznaranja]
    \textbf{\color{bordeadv}Atención:}\quad
}{%
    \end{mdframed}
}

\newenvironment{vizdef}{%
    \begin{mdframed}[style=vizazul]
}{%
    \end{mdframed}
}

% ── Comandos matemáticos ─────────────────────────────────────────────
\newcommand{\R}{\mathbb{R}}
\newcommand{\norm}[1]{\left\lVert #1 \right\rVert}
\newcommand{\grad}{\nabla}
\newcommand{\T}{{\top}}
\newcommand{\btheta}{\boldsymbol{\theta}}
\newcommand{\mstar}{m^{*}}
\newcommand{\bstar}{b^{*}}

% ── Hyperref ─────────────────────────────────────────────────────────
\hypersetup{
    colorlinks = true,
    linkcolor  = azulviz,
    urlcolor   = azulviz,
    citecolor  = azulviz
}

% ── Metadatos ────────────────────────────────────────────────────────
\title{
    {\Large\bfseries ¿Por qué el error cuadrático medio?}\\[0.4em]
    {\large Regresión lineal como máxima verosimilitud}\\[0.8em]
    {\normalsize Notas de lectura para \textit{GradienViz}}
}
\author{
    Departamento Académico de Sistemas Computacionales\\
    Universidad Autónoma de Baja California Sur
}
\date{2026}

% ─────────────────────────────────────────────────────────────────────
\begin{document}

\maketitle
\tableofcontents
\vspace{1em}

\begin{vizadv}
Este documento complementa los otros de la colección GradienViz.
Se asume familiaridad con la función de densidad normal (gaussiana),
con logaritmos y con la idea general del descenso por gradiente.
\end{vizadv}

% =============================================================================
\section{Una pregunta incómoda}
% =============================================================================

Los otros documentos de esta colección toman el error cuadrático medio
como punto de partida:
\[
    J(m, b) = \frac{1}{n}\sum_{i=1}^{n}(m x_i + b - y_i)^2.
\]
Pero no dicen por qué. ¿Por qué elevar al cuadrado?
¿Por qué no usar el error absoluto $|m x_i + b - y_i|$,
el error cúbico, o simplemente la suma de errores sin valor absoluto?

Todas esas funciones miden, en algún sentido, qué tan lejos está la
recta de los datos. Sin embargo, solo una de ellas tiene una justificación
estadística precisa bajo supuestos explícitos sobre cómo se generaron
los datos. Este documento desarrolla esa justificación.

La respuesta corta: \textbf{el error cuadrático medio es la función
de costo natural cuando el ruido en los datos es gaussiano}.
Minimizarlo equivale a maximizar la probabilidad de haber observado
exactamente los datos que tenemos.

% =============================================================================
\section{El modelo probabilista}
% =============================================================================

En los otros documentos, los datos $(x_i, y_i)$ simplemente existen.
Aquí los modelamos como el resultado de un proceso aleatorio.

\subsection{El supuesto de ruido gaussiano}

Asumimos que cada observación $y_i$ se generó así:
\begin{equation}\label{eq:modelo}
    y_i = m x_i + b + \varepsilon_i,
    \qquad \varepsilon_i \overset{\text{iid}}{\sim} \mathcal{N}(0, \sigma^2),
\end{equation}
donde $\varepsilon_i$ es un término de ruido aleatorio, independiente
e idénticamente distribuido según una distribución normal de media cero
y varianza $\sigma^2$.

Este supuesto tiene una interpretación física directa: la recta $y = mx + b$
describe la tendencia real, y $\varepsilon_i$ captura todo lo demás\textemdash
errores de medición, variables no incluidas en el modelo, variabilidad
inherente del fenómeno. Asumir que ese ruido es gaussiano, centrado en
cero y de varianza constante, es la hipótesis más común y más estudiada
en estadística.

\subsection{La verosimilitud de una observación}

Bajo el modelo~\eqref{eq:modelo}, dado un par de parámetros $(m, b)$,
cada observación $y_i$ sigue una distribución normal con media $mx_i + b$
y varianza $\sigma^2$:
\[
    y_i \mid x_i,\, m,\, b \;\sim\; \mathcal{N}(mx_i + b,\; \sigma^2).
\]
Su función de densidad es:
\begin{equation}\label{eq:densidad}
    f(y_i \mid x_i, m, b) = \frac{1}{\sigma\sqrt{2\pi}}
    \exp\!\left(-\frac{(y_i - mx_i - b)^2}{2\sigma^2}\right).
\end{equation}

\subsection{La función de verosimilitud}

La \textbf{verosimilitud} $L(m, b)$ mide qué tan probable es haber
observado \emph{todos} los datos $(x_1, y_1), \ldots, (x_n, y_n)$
si los parámetros verdaderos fueran $(m, b)$. Por independencia de
las observaciones, es el producto de las densidades individuales:
\begin{equation}\label{eq:likelihood}
    L(m, b) = \prod_{i=1}^{n} f(y_i \mid x_i, m, b)
    = \prod_{i=1}^{n} \frac{1}{\sigma\sqrt{2\pi}}
      \exp\!\left(-\frac{(y_i - mx_i - b)^2}{2\sigma^2}\right).
\end{equation}

El \textbf{estimador de máxima verosimilitud} (EMV) es el par $(m, b)$
que maximiza $L(m, b)$: los parámetros que hacen más probable haber
obtenido exactamente los datos observados.

\begin{viznota}{Panel~1}
Cada punto azul $(x_i, y_i)$ es una realización del modelo~\eqref{eq:modelo}.
La distancia vertical de cada punto a la recta naranja es el residuo
$e_i = y_i - (mx_i + b)$, que bajo el modelo es una realización de
$\varepsilon_i \sim \mathcal{N}(0, \sigma^2)$. Residuos grandes son
improbables bajo una gaussiana; el EMV los penaliza cuadráticamente.
\end{viznota}

% =============================================================================
\section{Maximizar \texorpdfstring{$L$}{L} equivale a minimizar MSE}
% =============================================================================

Trabajar directamente con el producto~\eqref{eq:likelihood} es
inconveniente: productos de exponenciales son difíciles de derivar.
La solución estándar es tomar logaritmo. Como el logaritmo es una
función estrictamente creciente, maximizar $L$ y maximizar
$\ell = \ln L$ tienen exactamente el mismo argumento óptimo.

\subsection{La log-verosimilitud}

Aplicando $\ln$ a~\eqref{eq:likelihood}:
\begin{align}
    \ell(m, b) &= \ln L(m, b)
    = \sum_{i=1}^{n} \ln f(y_i \mid x_i, m, b) \notag\\
    &= \sum_{i=1}^{n} \left[
        -\ln(\sigma\sqrt{2\pi})
        - \frac{(y_i - mx_i - b)^2}{2\sigma^2}
    \right] \notag\\
    &= -n\ln(\sigma\sqrt{2\pi})
       - \frac{1}{2\sigma^2}\sum_{i=1}^{n}(y_i - mx_i - b)^2.
    \label{eq:loglik}
\end{align}

\subsection{El resultado central}

En~\eqref{eq:loglik}, el primer término $-n\ln(\sigma\sqrt{2\pi})$
es una constante que no depende de $(m, b)$. El segundo término es
$-\frac{1}{2\sigma^2}\sum_{i=1}^n (y_i - mx_i - b)^2$, que tampoco
depende de $\sigma$ en cuanto a su argumento óptimo.

Por lo tanto:
\[
    \underbrace{\arg\max_{m,b}\; \ell(m,b)}_{\text{máxima verosimilitud}}
    \;=\;
    \arg\max_{m,b}\left[-\sum_{i=1}^{n}(y_i - mx_i - b)^2\right]
    \;=\;
    \underbrace{\arg\min_{m,b}\;\frac{1}{n}\sum_{i=1}^{n}(y_i - mx_i - b)^2}_{\text{mínimos cuadrados}}.
\]

\begin{vizdef}
\textbf{Equivalencia fundamental.} Bajo el supuesto de ruido gaussiano
independiente e idénticamente distribuido, el estimador de máxima
verosimilitud y el estimador de mínimos cuadrados son idénticos:
\[
    (\mstar, \bstar) = \arg\max_{m,b} L(m,b) = \arg\min_{m,b} J(m,b).
\]
Minimizar el error cuadrático medio no es una elección arbitraria:
es la consecuencia directa de asumir ruido gaussiano.
\end{vizdef}

\begin{viznota}{Panel~2, punto $\oplus$}
La cruz amarilla marca $(\mstar, \bstar)$: el mínimo de $J(m,b)$.
Por la equivalencia anterior, ese mismo punto es también el máximo
de la verosimilitud $L(m,b)$ bajo el modelo gaussiano. Es el punto
que hace más probable haber observado los datos del Panel~1.
\end{viznota}

% =============================================================================
\section{Cambia el ruido, cambia la función de costo}
% =============================================================================

La derivación anterior revela algo más profundo: el MSE no es
universal. Es la función de costo correcta \emph{para el ruido gaussiano}.
Si el mecanismo que genera los errores es distinto, la función de costo
que maximiza la verosimilitud también lo es.

\subsection{Ruido de Laplace: el error absoluto medio}

Si en lugar de $\varepsilon_i \sim \mathcal{N}(0, \sigma^2)$ asumimos
$\varepsilon_i \sim \text{Laplace}(0, b)$, cuya densidad es
$f(\varepsilon) = \frac{1}{2b}e^{-|\varepsilon|/b}$, entonces la
log-verosimilitud resulta proporcional a $-\sum_i |y_i - mx_i - b|$.
El estimador de máxima verosimilitud minimiza el
\textbf{error absoluto medio} (MAE):
\[
    J_{\text{MAE}}(m, b) = \frac{1}{n}\sum_{i=1}^{n}|y_i - mx_i - b|.
\]
El MAE es más robusto ante valores atípicos (\emph{outliers}) que el MSE,
precisamente porque la distribución de Laplace asigna mayor probabilidad
a errores grandes que la gaussiana.

\subsection{Modelo de clasificación binaria: la entropía cruzada}

En clasificación, $y_i \in \{0, 1\}$ y el modelo predice una
probabilidad $\hat{p}_i = \sigma(\btheta^\T \mathbf{x}_i) \in (0,1)$.
Cada $y_i$ sigue una distribución de Bernoulli con parámetro $\hat{p}_i$.
La log-verosimilitud negativa de este modelo es la
\textbf{entropía cruzada binaria}:
\[
    J_{\text{CE}}(\btheta) = -\frac{1}{n}\sum_{i=1}^{n}
    \bigl[y_i \ln \hat{p}_i + (1 - y_i)\ln(1 - \hat{p}_i)\bigr].
\]
Esta es exactamente la función de costo que minimiza el descenso por
gradiente en regresión logística y en la capa de salida de una red
neuronal para clasificación.

\subsection{El origen estadístico de las funciones de costo}

\begin{center}
\small
\begin{tabular}{llll}
    \toprule
    \textbf{Supuesto de ruido} & \textbf{Distribución} &
    \textbf{Función de costo} & \textbf{Uso típico} \\
    \midrule
    Gaussiano   & $\mathcal{N}(0,\sigma^2)$   & MSE (error cuadrático)  & Regresión \\
    Laplace     & $\text{Laplace}(0,b)$       & MAE (error absoluto)    & Regresión robusta \\
    Bernoulli   & $\text{Bernoulli}(\hat{p})$ & Entropía cruzada binaria & Clasificación binaria \\
    Categórica  & $\text{Categorical}(\hat{p})$ & Entropía cruzada       & Clasificación múltiple \\
    \bottomrule
\end{tabular}
\end{center}

\begin{vizadv}
En la práctica, elegir una función de costo equivale a elegir un
supuesto sobre la distribución del ruido. Cuando un modelo de
aprendizaje automático usa entropía cruzada en lugar de MSE, no es
por conveniencia computacional: es porque el problema tiene estructura
de clasificación (ruido de Bernoulli) en lugar de regresión (ruido
gaussiano). La función de costo codifica una hipótesis estadística
sobre los datos.
\end{vizadv}

% =============================================================================
\section*{Resumen}
% =============================================================================

El error cuadrático medio no es una elección arbitraria.
Es la función de costo que emerge naturalmente cuando se asume que
los errores del modelo siguen una distribución gaussiana centrada en
cero: minimizar MSE equivale a encontrar los parámetros que maximizan
la verosimilitud de los datos observados bajo ese supuesto.

La cruz $\oplus$ que GradienViz marca en el Panel~2 tiene, por lo
tanto, dos lecturas simultáneas y equivalentes:
el mínimo de la función de costo $J(m,b)$ y
el máximo de la verosimilitud $L(m,b)$.

Cambiar el supuesto de ruido cambia la función de costo, pero no
cambia el método: el descenso por gradiente sigue siendo el algoritmo
que minimiza cualquiera de ellas.

\end{document}
